\documentclass[12pt]{article}
\usepackage[margin=1in]{geometry}
\usepackage{setspace}
\usepackage[T1]{fontenc}
\usepackage{newtxtext,newtxmath}
\usepackage{xcolor}
\usepackage{colortbl}
\usepackage{enumitem}
\usepackage{hyperref}
\doublespacing
\setlength{\parindent}{0pt}
\setlength{\parskip}{0.2em}
\setlistdepth{6}
\renewlist{itemize}{itemize}{6}
\setlist[itemize,1]{label=\textbullet,leftmargin=1.6em,itemsep=0.15em,topsep=0.15em}
\setlist[itemize,2]{label=\textbullet,leftmargin=1.9em,itemsep=0.15em,topsep=0.15em}
\setlist[itemize,3]{label=\textbullet,leftmargin=2.2em,itemsep=0.15em,topsep=0.15em}
\setlist[itemize,4]{label=\textbullet,leftmargin=2.5em,itemsep=0.15em,topsep=0.15em}
\setlist[itemize,5]{label=\textbullet,leftmargin=2.8em,itemsep=0.15em,topsep=0.15em}
\setlist[itemize,6]{label=\textbullet,leftmargin=3.1em,itemsep=0.15em,topsep=0.15em}
\definecolor{nikhicolor}{HTML}{8E7CC3}
\definecolor{samcolor}{HTML}{6AA84F}
\definecolor{pitercolor}{HTML}{CC6666}
\begin{document}
\begin{center}
{\Large \textbf{HTSA Group Paper Outline}}\\
{\normalsize BIOL550 Group Project (Draft 1)}\\[0.5em]
{\normalsize Nikhi Boggavarapu}\\
{\normalsize Sam Kopelev}\\
{\normalsize Piter Garcia}\\
\end{center}
\vspace{0.5em}
\noindent
\begin{tabular}{|p{1.55in}|p{4.65in}|}
\hline
\textbf{Member} & \textbf{Note} \\ \hline
\cellcolor[HTML]{D9D2E9}Nikhi Boggavarapu & \cellcolor[HTML]{D9D2E9}Highlighted text in light purple indicates Nikhi's contribution. \\ \hline
\cellcolor[HTML]{B6D7A8}Sam Kopelev & \cellcolor[HTML]{B6D7A8}Highlighted text in light green indicates Sam's contribution. \\ \hline
\cellcolor[HTML]{F4CCCC}Piter Garcia & \cellcolor[HTML]{F4CCCC}Highlighted text in light red indicates Piter's contribution. \\ \hline
\end{tabular}
\section*{\textcolor{nikhicolor}{Introduction}}
\begin{itemize}
\item \textcolor{nikhicolor}{Brief intro about high-throughput sequencing}
\begin{itemize}
\item \textcolor{pitercolor}{RNA-seq as a transcriptome-wide method for detecting gene expression changes}
\item \textcolor{nikhicolor}{Types of NGS}
\begin{itemize}
\item \textcolor{nikhicolor}{Organized by read length}
\item \textcolor{nikhicolor}{Applications}
\end{itemize}
\item \textcolor{nikhicolor}{Benefits of NGS}
\begin{itemize}
\item \textcolor{nikhicolor}{High throughput}
\item \textcolor{nikhicolor}{Fast}
\end{itemize}
\item \textcolor{nikhicolor}{What our goal was in utilizing NGS}
\begin{itemize}
\item \textcolor{pitercolor}{Unbiased, transcriptome-wide detection of gene expression changes}
\end{itemize}
\end{itemize}
\item \textcolor{samcolor}{Introduce paper/dataset}
\begin{itemize}
\item \textcolor{samcolor}{Experimental design}
\item \textcolor{samcolor}{DRG after spinal cord injury}
\item \textcolor{samcolor}{Goals of the paper}
\begin{itemize}
\item \textcolor{samcolor}{Why did they analyze DRG with cKO}
\item \textcolor{samcolor}{Ipsilateral vs contralateral}
\item \textcolor{pitercolor}{FF / cre context}
\end{itemize}
\end{itemize}
\item \textcolor{samcolor}{How we wanted to use the paper}
\begin{itemize}
\item \textcolor{samcolor}{Explore and confirm suggested pathways}
\begin{itemize}
\item \textcolor{samcolor}{FF / cre}
\end{itemize}
\item \textcolor{samcolor}{Examine gene expression changes beyond the genes emphasized in the original paper}
\end{itemize}
\item \textcolor{pitercolor}{Claim}
\begin{itemize}
\item \textcolor{pitercolor}{\textit{Differential expression (DE):} transcriptome-wide RNA-seq helps identify genes involved beyond the ones the paper discussed}
\item \textcolor{pitercolor}{\textit{Key contrasts:} expression differences were quantified across the main experimental comparisons, especially ipsilateral vs contralateral DRG, and genotype where relevant}
\item \textcolor{pitercolor}{\textit{GO / pathway enrichment:} DE gene sets were interpreted in terms of broader biological processes and pathways linked to injury response and the cKO context}
\end{itemize}
\end{itemize}
\section*{\textcolor{pitercolor}{Materials and Methods}}
\begin{itemize}
\item \textcolor{pitercolor}{Dataset and study design}
\begin{itemize}
\item \textcolor{pitercolor}{Public dataset accession}
\item \textcolor{pitercolor}{Tissue and injury context}
\item \textcolor{pitercolor}{Sample groups and contrasts}
\item \textcolor{pitercolor}{Why was this subset retained}
\end{itemize}
\item \textcolor{nikhicolor}{Pipeline Steps}
\begin{itemize}
\item \textcolor{nikhicolor}{Data Collection and Preprocessing}
\begin{itemize}
\item \textcolor{nikhicolor}{Download SRR files (`.gz`)}
\item \textcolor{pitercolor}{FASTQ organization}
\item \textcolor{pitercolor}{Sample manifest / metadata table}
\item \textcolor{pitercolor}{Tools used}
\end{itemize}
\item \textcolor{nikhicolor}{Data Quality Control and Trimming}
\begin{itemize}
\item \textcolor{nikhicolor}{Raw FastQC / MultiQC}
\item \textcolor{nikhicolor}{FASTP trimming}
\item \textcolor{nikhicolor}{Post-trim FastQC / MultiQC}
\item \textcolor{nikhicolor}{Key QC checkpoint metrics}
\end{itemize}
\item \textcolor{nikhicolor}{Data Preparation (Alignment and count generation)}
\begin{itemize}
\item \textcolor{nikhicolor}{Reference genome and annotation}
\begin{itemize}
\item \textcolor{nikhicolor}{(mm39)}
\end{itemize}
\item \textcolor{nikhicolor}{STAR index and alignment}
\begin{itemize}
\item \textcolor{nikhicolor}{Building the STAR genome}
\item \textcolor{nikhicolor}{Paired-end alignment}
\end{itemize}
\item \textcolor{pitercolor}{GeneCounts output}
\item \textcolor{pitercolor}{BAM / log / count outputs}
\item \textcolor{pitercolor}{Family count matrix assembly}
\end{itemize}
\item \textcolor{samcolor}{Data Analysis and Interpretation}
\begin{itemize}
\item \textcolor{samcolor}{MultiQC}
\item \textcolor{samcolor}{Differential Expression Analysis}
\begin{itemize}
\item \textcolor{samcolor}{DESeq2 setup}
\item \textcolor{samcolor}{DESeq2}
\begin{itemize}
\item \textcolor{samcolor}{Volcano plot}
\item \textcolor{samcolor}{Heatmap of top DEGs (Regeneration-Enhancing)}
\item \textcolor{pitercolor}{MA plot}
\item \textcolor{samcolor}{Target validation (gene expression boxplots)}
\end{itemize}
\item \textcolor{pitercolor}{Filtering rule}
\item \textcolor{pitercolor}{Modeled contrasts}
\item \textcolor{pitercolor}{Thresholds / significance criteria}
\end{itemize}
\item \textcolor{samcolor}{Follow-up selection for functional interpretation}
\begin{itemize}
\item \textcolor{pitercolor}{Bend-point rule}
\item \textcolor{samcolor}{GO analysis}
\item \textcolor{samcolor}{GSEA analysis}
\item \textcolor{samcolor}{Panther analysis}
\item \textcolor{samcolor}{Extra analysis, if finalized}
\end{itemize}
\end{itemize}
\end{itemize}
\end{itemize}
\section*{\textcolor{samcolor}{Results}}
\begin{itemize}
\item \textcolor{samcolor}{Dataset quality supported downstream analysis}
\begin{itemize}
\item \textcolor{samcolor}{Number of samples and read structure}
\item \textcolor{samcolor}{Pre-trim}
\begin{itemize}
\item \textcolor{samcolor}{GC content}
\item \textcolor{samcolor}{Adapter sequences}
\end{itemize}
\item \textcolor{samcolor}{Alignment stats}
\begin{itemize}
\item \textcolor{samcolor}{Unique mapping}
\end{itemize}
\end{itemize}
\item \textcolor{pitercolor}{The sample structure showed the strongest separation by the main biological contrast}
\begin{itemize}
\item \textcolor{pitercolor}{PCA (key figure)}
\item \textcolor{pitercolor}{Interpretation of PC1 / PC2}
\item \textcolor{pitercolor}{Genotype analysis}
\end{itemize}
\item \textcolor{samcolor}{Differential expression identified the strongest transcriptomic changes}
\begin{itemize}
\item \textcolor{samcolor}{PCA (key figure)}
\begin{itemize}
\item \textcolor{samcolor}{Explain how it is separated, PC1 and PC2}
\end{itemize}
\item \textcolor{nikhicolor}{MA plot}
\item \textcolor{nikhicolor}{Cumulative Distribution Plot (bend-point / elbow rule)}
\item \textcolor{nikhicolor}{Volcano Plot (Important Genes)}
\begin{itemize}
\item \textcolor{nikhicolor}{Include a secondary Volcano Plot with a new threshold}
\end{itemize}
\item \textcolor{nikhicolor}{Heatmap}
\begin{itemize}
\item \textcolor{nikhicolor}{Including the interpretation of the distance heatmap between the two conditions}
\end{itemize}
\end{itemize}
\item \textcolor{pitercolor}{Bend-point selection narrowed the main follow-up sets}
\begin{itemize}
\item \textcolor{pitercolor}{Cumulative plot}
\item \textcolor{pitercolor}{Selected gene counts}
\item \textcolor{pitercolor}{Why this helped interpretation}
\item \textcolor{pitercolor}{Secondary volcano plot, if you keep it}
\end{itemize}
\item \textcolor{nikhicolor}{Functional enrichment connected DE results to broader biology}
\begin{itemize}
\item \textcolor{nikhicolor}{GO / pathway results}
\item \textcolor{nikhicolor}{Pathway level / Panther analysis}
\begin{itemize}
\item \textcolor{nikhicolor}{Go through the plots listed by Panther}
\end{itemize}
\item \textcolor{nikhicolor}{Proteostasis}
\item \textcolor{nikhicolor}{Translation}
\item \textcolor{nikhicolor}{Metabolism}
\item \textcolor{nikhicolor}{Extra analysis, once finalized}
\end{itemize}
\end{itemize}
\section*{\textcolor{nikhicolor}{Discussion}}
\begin{itemize}
\item \textcolor{pitercolor}{Main biological interpretation (what does the data show)}
\begin{itemize}
\item \textcolor{pitercolor}{What the data show overall}
\item \textcolor{pitercolor}{Strongest supported signal}
\item \textcolor{pitercolor}{What is primary vs secondary}
\item \textcolor{pitercolor}{Gene expression differences}
\begin{itemize}
\item \textcolor{pitercolor}{Injury vs control}
\end{itemize}
\item \textcolor{pitercolor}{Pathways being affected}
\begin{itemize}
\item \textcolor{pitercolor}{Proteostasis (AhR activation)}
\item \textcolor{pitercolor}{Translation (suppression / upregulation of genes)}
\item \textcolor{pitercolor}{Metabolism (energy output)}
\end{itemize}
\end{itemize}
\item \textcolor{nikhicolor}{What was unusual or needs caution (things that were weird about the dataset to consider)}
\begin{itemize}
\item \textcolor{nikhicolor}{Odd volcano plots}
\item \textcolor{nikhicolor}{The genotype signal appears weaker / secondary}
\item \textcolor{nikhicolor}{PCA collisions / overlap}
\item \textcolor{nikhicolor}{Interpretation vs causation}
\end{itemize}
\item \textcolor{pitercolor}{What this adds beyond the original paper}
\begin{itemize}
\item \textcolor{pitercolor}{Expands the analysis beyond the paper’s candidate genes}
\begin{itemize}
\item \textcolor{samcolor}{Helps link pathways}
\begin{itemize}
\item \textcolor{samcolor}{Proposed and new}
\end{itemize}
\end{itemize}
\item \textcolor{pitercolor}{Helps link the pathways already proposed in the paper with additional transcriptomic signals}
\item \textcolor{pitercolor}{Connects newly identified genes to the genes highlighted in the paper through broader transcriptomic patterns}
\item \textcolor{pitercolor}{Moves interpretation from single-gene emphasis toward pathway- and network-level context}
\end{itemize}
\item \textcolor{samcolor}{NGS in the context of this paper}
\begin{itemize}
\item \textcolor{samcolor}{Global discovery}
\item \textcolor{samcolor}{Reproducibility}
\item \textcolor{samcolor}{External applications / biological relevance}
\begin{itemize}
\item \textcolor{samcolor}{How does our global discovery connect to the genes found in the paper}
\item \textcolor{samcolor}{Two-hybrid screening}
\end{itemize}
\end{itemize}
\item \textcolor{pitercolor}{Future validation (what it suggests for future validation)}
\begin{itemize}
\item \textcolor{pitercolor}{qPCR for key differentially expressed genes}
\item \textcolor{pitercolor}{Functional follow-up experiments for injury-response candidates}
\item \textcolor{pitercolor}{Two-hybrid screening only if it remains biologically justified}
\end{itemize}
\item \textcolor{pitercolor}{Limitations and cautions}
\begin{itemize}
\item \textcolor{pitercolor}{Weaker genotype signal}
\item \textcolor{pitercolor}{PCA overlap / collisions}
\item \textcolor{pitercolor}{Dependence on the original dataset design}
\item \textcolor{pitercolor}{Interpretation vs causation}
\end{itemize}
\end{itemize}
\section*{References (APA)}
\begin{itemize}
\item \href{https://www.nature.com/articles/s41586-026-10295-z\#code-availability}{Main Paper}
\item NGS overview / background sources
\item Additional papers using or reviewing NGS
\item Background sources for the mouse study
\item Minimum of 15 primary sources in the final paper
\end{itemize}
\section*{Code Availability}
\begin{itemize}
\item Git repository
\item Notebook(s)
\end{itemize}
\section*{Source Data}
\begin{itemize}
\item CSVs / tables supporting plots
\item Data frames / exported tables
\end{itemize}
\section*{Supplemental Material}
\begin{itemize}
\item Full figure set
\item Additional supporting outputs
\end{itemize}
\end{document}
